\documentclass[11pt,a4paper]{article}
\usepackage[T1]{fontenc}
\usepackage[utf8]{inputenc}
\usepackage{amsmath,amssymb,amsthm}
\usepackage[margin=1in]{geometry}
\usepackage[colorlinks=true,linkcolor=blue,urlcolor=blue]{hyperref}
\theoremstyle{plain}
\newtheorem{theorem}{Theorem}
\newtheorem{lemma}[theorem]{Lemma}
\newtheorem{proposition}[theorem]{Proposition}
\newtheorem{corollary}[theorem]{Corollary}
\theoremstyle{definition}
\newtheorem{remark}[theorem]{Remark}

\title{The empirical recurrence for OEIS A252641, proved by transfer matrix}
\author{Adrian Perez Fontelles\\ \small Independent researcher}
\date{1 September 2026}
\begin{document}
\maketitle

\begin{abstract}
OEIS A252641 counts $3\times(n+2)$ arrays over $\{0,\dots,3\}$ subject to a condition imposed on
every $3\times3$ subblock, and it carries an empirical recurrence of order $67$
contributed by R.~H.~Hardin. It is true, and it is not an empirical matter. A $3\times3$
subblock spans three consecutive columns, so the state of a transfer matrix is a PAIR of
consecutive columns and appending one more column is a step; an array is then a walk, and
$a(n)$ is a walk count on $S=4096$ vertices. Such a sequence is $C$-finite, so whether a given
recurrence annihilates it is decided by a finite exact computation, with the Cayley--Hamilton
theorem bounding the work.
\end{abstract}

\noindent\small 2020 Mathematics Subject Classification. 05A15, 05B45, 15B36.\normalsize

\section{The conjecture}

OEIS A252641 is ``Number of (1+2) X (n+2) 0..3 arrays with every 3 X 3 subblock row and diagonal sum equal to 0 3 4 6 or 7 and every 3 X 3 column and antidiagonal sum not equal to 0 3 4 6 or 7.''. It has offset $1$ and begins
\[
572, 1013, 1758, 3162, 6598, 13038, 24852, 51396,\ \dots
\]
The entry states, as a conjecture contributed by its author and never marked settled:
\begin{quote}
Empirical: a(n) = 54*a(n-3) -1300*a(n-6) +a(n-7) +a(n-8) +18442*a(n-9) -46*a(n-10) -48*a(n-11) -171869*a(n-12) +939*a(n-13) +1007*a(n-14) +1112232*a(n-15) -11261*a(n-16) -12125*a(n-17) -5161835*a(n-18) +88673*a(n-19) +92394*a(n-20) +17537697*a(n-21) -486813*a(n-22) -461456*a(n-23) -44251113*a(n-24) +1941080*a(n-25) +1492371*a(n-26) +83816227*a(n-27) -5824359*a(n-28) -2806042*a(n-29) -120301341*a(n-30) +13608098*a(n-31) +1230822*a(n-32) +132500058*a(n-33) -25409661*a(n-34) +8866119*a(n-35) -115158185*a(n-36) +38073643*a(n-37) -28401518*a(n-38) +84586799*a(n-39) -44658830*a(n-40) +46554872*a(n-41) -59077105*a(n-42) +39166599*a(n-43) -48767596*a(n-44) +41880073*a(n-45) -24271783*a(n-46) +33857442*a(n-47) -27389506*a(n-48) +9919627*a(n-49) -15168789*a(n-50) +13979868*a(n-51) -2423220*a(n-52) +4029102*a(n-53) -4453677*a(n-54) +361068*a(n-55) -512244*a(n-56) +193374*a(n-57) -111429*a(n-58) -17496*a(n-59) +514188*a(n-60) -1998*a(n-61) +23328*a(n-62) -196344*a(n-63) +61236*a(n-64) +23328*a(n-66) -21384*a(n-67) for n>74.
\end{quote}
As of the ``Last modified'' line on the live entry (Nov 03 2022 21:49:39, revision 7) this is still
recorded as empirical, and nothing on the entry records it as proved.

\section{Three consecutive lines}

Write an admissible array as a sequence of columns $u_1,u_2,\dots$, each an element of
$\{0,\dots,3\}^{3}$. A $3\times3$ subblock occupies three consecutive columns and
three consecutive rows, so with $r,s,t$ three consecutive columns the blocks it
contributes are
\[
g=\begin{pmatrix}r_i&s_i&t_i\\ r_{i+1}&s_{i+1}&t_{i+1}\\ r_{i+2}&s_{i+2}&t_{i+2}\end{pmatrix}
\]
for each admissible index, and the entry's condition is the requirement that
\[
\text{every row, diagonal sum} \in \{0,3,4,6,7\}\ \text{ and }\ \text{every column, antidiagonal sum} \notin \{0,3,4,6,7\}
\]
holds for every such $g$. This involves $r$, $s$ and $t$ only: it is a condition on
three consecutive columns and on nothing else.

Take as vertices of a digraph $G$ the ordered PAIRS $(r,s)$ of columns, of which there are
$S=(3+1)^{2\cdot3}=4096$, and put an edge $(r,s)\to(s,t)$ exactly when the condition
holds for the triple $(r,s,t)$ at every admissible index. An array with $L$ columns is then a
walk of length $L-2$ in $G$: its first two columns name the starting vertex and each further
column is one step. Hence, with $M$ the adjacency matrix of $G$,
\[
a(n)\;=\;\,\mathbf{1}^{\!\top}M^{\,n-1+1}\mathbf{1},
\]
the exponent being fixed by the entry's own shape and checked against its published terms.
In particular $a$ is $C$-finite.

\section{The proof}

\begin{lemma}\label{lem:crit}
Let $q(t)=t^{67}-\sum_i c_it^{67-i}$ be the characteristic polynomial of the
conjectured recurrence. Then the residual $a(n)-\sum_ic_ia(n-i)$ equals
$\mathbf{1}^{\!\top}M^{\,j}q(M)\mathbf{1}$ for the corresponding walk index
$j$, so the recurrence holds from some point on if and only if
$\mathbf{1}^{\!\top}M^{j}q(M)\mathbf{1}=0$ for all large $j$; and by the
Cayley--Hamilton theorem it is enough to examine $j<S$.
\end{lemma}

\begin{proof}
Factor $M^{j}$ out of $M^{j+67}-\sum_ic_iM^{j+67-i}$. For the bound, $M^{S}$
is an integer combination of $I,M,\dots,M^{S-1}$, so the numbers
$u_j=\mathbf{1}^{\!\top}M^{j}q(M)\mathbf{1}$ satisfy the monic linear recurrence given by
the characteristic polynomial of $M$; $S$ consecutive zeros therefore force every later one,
and the last nonzero $u_j$ pins the threshold exactly.
\end{proof}

\begin{theorem}
\[
a(n)\;=\;54\,a(n-3) - 1300\,a(n-6) + a(n-7) + a(n-8) + 18442\,a(n-9) - 46\,a(n-10) - 48\,a(n-11) - 171869\,a(n-12) + 939\,a(n-13) + 1007\,a(n-14) + 1112232\,a(n-15) - 11261\,a(n-16) - 12125\,a(n-17) - 5161835\,a(n-18) + 88673\,a(n-19) + 92394\,a(n-20) + 17537697\,a(n-21) - 486813\,a(n-22) - 461456\,a(n-23) - 44251113\,a(n-24) + 1941080\,a(n-25) + 1492371\,a(n-26) + 83816227\,a(n-27) - 5824359\,a(n-28) - 2806042\,a(n-29) - 120301341\,a(n-30) + 13608098\,a(n-31) + 1230822\,a(n-32) + 132500058\,a(n-33) - 25409661\,a(n-34) + 8866119\,a(n-35) - 115158185\,a(n-36) + 38073643\,a(n-37) - 28401518\,a(n-38) + 84586799\,a(n-39) - 44658830\,a(n-40) + 46554872\,a(n-41) - 59077105\,a(n-42) + 39166599\,a(n-43) - 48767596\,a(n-44) + 41880073\,a(n-45) - 24271783\,a(n-46) + 33857442\,a(n-47) - 27389506\,a(n-48) + 9919627\,a(n-49) - 15168789\,a(n-50) + 13979868\,a(n-51) - 2423220\,a(n-52) + 4029102\,a(n-53) - 4453677\,a(n-54) + 361068\,a(n-55) - 512244\,a(n-56) + 193374\,a(n-57) - 111429\,a(n-58) - 17496\,a(n-59) + 514188\,a(n-60) - 1998\,a(n-61) + 23328\,a(n-62) - 196344\,a(n-63) + 61236\,a(n-64) + 23328\,a(n-66) - 21384\,a(n-67)
\]
for every $n>74$.
\end{theorem}

\begin{proof}
The vector $w=q(M)\mathbf{1}$ was formed by $67$ matrix--vector products in exact
integer arithmetic and $\mathbf{1}^{\!\top}M^{j}w$ evaluated for $j=0,1,\dots$, again
exactly. Those integers vanish from the index corresponding to $n=74+1$ onwards, and the
run was continued until the working vector was identically zero, which settles all larger $j$
at once. By Lemma~\ref{lem:crit} the recurrence holds for every $n>74$.
\end{proof}

\section{Verification}

Three checks, each independent of the algebra above.

First, the transfer model reproduces all $26$ terms the entry publishes, exactly, in
integer arithmetic. This is what ties the model to the entry: the digraph is built by reading
the entry's English, and a misreading gives different counts.

Second, the conjectured recurrence was evaluated directly on the entry's own published terms,
with no matrices involved, and holds wherever the proved range and the data overlap.

Third, the annihilation test of Lemma~\ref{lem:crit} was carried out in exact integer
arithmetic on the whole vector rather than on a sampled prefix, so no rounding or sampling
enters anywhere.

\begin{thebibliography}{9}
\bibitem{oeis} The OEIS Foundation, \emph{The On-Line Encyclopedia of Integer
Sequences}, \url{https://oeis.org}, 2026. Sequence A252641.
\bibitem{stanley} R.~P.~Stanley, \emph{Enumerative Combinatorics}, Volume 1, 2nd ed.,
Cambridge University Press, 2011. (Transfer-matrix method, Section 4.7.)
\bibitem{fs} P.~Flajolet and R.~Sedgewick, \emph{Analytic Combinatorics}, Cambridge
University Press, 2009.
\bibitem{horn} R.~A.~Horn and C.~R.~Johnson, \emph{Matrix Analysis}, 2nd ed., Cambridge
University Press, 2013.
\end{thebibliography}

\end{document}
